% Part: computability
% Chapter: computability-theory
% Section: complement-ce

\documentclass[../../../include/open-logic-section]{subfiles}

\begin{document}

\olfileid{cmp}{thy}{cmp}
\olsection{Himpunan kang Bisa Dienumerasi sacara Komputabel Ora Katutup tumrap Komplemen}

Upamane $A$ bisa dienumerasi sacara komputabel.
Apa komplemen~$A$,
$\Complement{A} = \Nat \setminus A$,
uga mesthi bisa dienumerasi sacara komputabel?
Teorema lan akibat ing ngisor iki nuduhake
yen wangsulane ``ora''.

\begin{thm}
\ollabel{thm:ce-comp}
Tetepna $A$ minangka himpunan sembarang
saka wilangan asli. Mula $A$ komputabel yen
lan mung yen $A$ lan $\Complement{A}$
loro-lorone bisa dienumerasi sacara komputabel.
\end{thm}

\begin{proof}
Arah maju gampang: yen $A$ komputabel,
mula $\Complement{A}$ uga komputabel
($\Char{A} = 1 \tsub
\Char{\Complement{A}}$), saengga loro-lorone
bisa dienumerasi sacara komputabel. Cathetan panyunting OLPL-766: persamaan kang dicithak ngasilake karakteristik A saka karakteristik komplemen; kanggo arah bukti iki, tukarna perane, saengga karakteristik komplemen diitung minangka siji dikurangi karakteristik A.

Ing arah kosok baline, upamane $A$ lan
$\Complement{A}$ loro-lorone bisa
dienumerasi sacara komputabel. Tetepna
$A$ minangka domain~$\cfind{d}$, lan
$\Complement{A}$ minangka domain~$\cfind{e}$.
Tetepna $h$ minangka
\[
h(x) = \umin{s}{(T(d,x,s) \lor T(e,x,s))}.
\]
Kanthi tembung liya, ing input~$x$,
$h$~nggoleki komputasi kang mandheg saka
$\cfind{d}$ utawa $\cfind{e}$. Yen
$x \in A$, panggolekan iki kasil ing
kasus kapisan; yen $x \in
\Complement{A}$, kasil ing kasus kapindho.
Mula $h$~fungsi komputabel total. Cathetan koreksi sumber OLPL-139: tes anggota A nggunakake indeks d, dudu indeks e kang ngenali komplemene.
% OLPL-139: Sumber beku nulis e ing tes anggota A, padahal d ngindeks A lan e ngindeks komplemene.
Nanging saiki kanggo saben~$x$, $x \in A$
yen lan mung yen $T(d, x, h(x))$,
yaiku yen $\cfind{d}$ kang ditegesi.
Amarga $T(d, x, h(x))$ relasi komputabel,
$A$~komputabel.
\end{proof}

\begin{explain}
Ing basa komputasi kang ora formal,
iki luwih gampang dimangerteni. Cathetan koreksi sumber OLPL-140: indeks ing andharan iki d kanggo A lan e kanggo komplemene, dudu e lan f kaya ing sumber. Kanggo
mutusake $A$, ing input $x$ goleka
% OLPL-140: Sumber beku nyebut e lan f, banjur ngarani e anggota A; indeks sing bener yaiku d kanggo A lan e kanggo komplemene.
komputasi kang mandheg saka $\cfind{d}$
lan $\cfind{e}$. Salah sijine mesthi
mandheg; yen iku $\cfind{d}$, mula
$x$ kalebu~$A$, yen ora, $x$ kalebu
$\Complement{A}$.
\end{explain}

\begin{cor}
\ollabel{cor:comp-k}
$\Complement{K_0}$ ora bisa dienumerasi
sacara komputabel.
\end{cor}

\begin{proof}
Kita ngerti yen $K_0$ bisa dienumerasi
sacara komputabel, nanging ora komputabel.
Yen $\Complement{K_0}$ bisa dienumerasi
sacara komputabel, mula $K_0$ mesthi
komputabel miturut \olref{thm:ce-comp},
nalisir \olref[ncp]{thm:K-0}.
\end{proof}

\end{document}
